\problemname{Plankton}
The fishes Alexa and Boris are going to start a plankton farm at the bottom of the sea. To prevent
the plankton from escaping, they want to build rectangular enclosures surrounded by fences.
Each enclosure must have an area of at least $A \text{ cm}^2$ so that the plankton has enough space to grow.

The fence for a given enclosure is set up along its entire perimeter in the form of four pre-built pieces:
a pair of equally long horizontal pieces and a pair of equally long vertical pieces. Alexa and Boris have
already purchased several pairs of fence pieces, where the pieces in each pair are of equal length.
They now want to build as many enclosures with an area of at least $A \text{ cm}^2$ as possible. Can
you help them calculate how many enclosures they can build simultaneously with their fence pieces?


\section*{Input}
The first line of input contains the integer $A$ ($1 \leq A \leq 10^9$), the minimum area that each enclosure must have.

The next line contains an integer $N$ ($1 \leq N \leq 100,000$), the number of pairs of fence pieces that
Alexa and Boris have already purchased.

The last line contains $N$ integers $l_1, l_2, \dots, l_N$ ($1 \leq l_i \leq 10^9$), where $l_i$ is the length
of the pieces in the $i$-th pair they bought, measured in centimeters.

\section*{Output}
Print an integer: the largest number of enclosures that Alexa and Boris can build using their fence pieces.

\section*{Points}
Your solution will be tested on several test case groups.
To get the points for a group, it must pass all the test cases in the group.

\noindent
\begin{tabular}{| l | l | p{12cm} |}
  \hline
  \textbf{Group} & \textbf{Point value} & \textbf{Constraints} \\ \hline
  $1$    & $20$       & The group consists of a single subtask, the one on our poster (\url{https://www.progolymp.se/2023/affisch.pdf}). \\ \hline
  $2$    & $12$       & $N \le 300$ \\ \hline
  $3$    & $28$       & $N \le 2000$ \\ \hline
  $4$    & $29$       & $A \le 2000$ \\ \hline
  $5$    & $11$       & No additional constraints. \\ \hline
\end{tabular}

\section*{Explanation of sample}
In the first example case, a maximum of 2 enclosures can be built.
For example, one can choose the pairs of fence pieces with lengths of $5$ cm and $13$ cm
(with an area of $5 \text{ cm} \times 13 \text{ cm} = 65 \text{ cm}^2$) and the pairs with
lengths of $8$ cm and $21$ cm (with an area of $8 \text{ cm} \times 21 \text{ cm} = 168 \text{ cm}^2$),
both of which are at least $40 \text{ cm}^2$.

In the second example case, all pairs of fence pieces have the same length.
However, with this length, the enclosure has an area of $100 \text{ cm}^2$, which is just slightly too small.

In the third case, Alexa and Boris are satisfied with an area of $100 \text{ cm}^2$, so all pairs of fence
pieces can be used to build a total of $5$ enclosures.
